\section{近期研究（2026 年 8--11 月）}
\paragraph{NGS-Marker：从局部高斯子集解码资产水印}
2026.08~\cite{2608_17447}。图像水印往往依赖特定渲染结果，资产只被取走一部分时未必仍可解码。NGS-Marker 学习高斯空间的消息注入与提取，并以预训练解码器监督随机局部片段的渐进嵌入，使被截取的高斯也携带可恢复的身份信息。渲染误差和消息恢复共同约束训练。

作者用 24 个训练场景和四个留出场景，在局部子集被提取并插入另一场景的设置下比较位准确率与原语级识别。方法需要访问高斯原语，超大场景高效分区与遍历仍是作者指出的扩展问题。

\paragraph{TRACE：用射线汇总量做跨私有子地图选视}
2026.10~\cite{2610_00822}。多机器人在各自持有高斯子地图时，某条候选射线的遮挡仍受其他机器人地图影响。TRACE 按深度区间交换光学厚度、辐射与位姿导数的汇总量，用来近似合并地图上的选视信息增益，无需发送整个原语集合。

在 Habitat-Sim 的 Gibson 室内场景，作者固定建图和规划管线，只比较选视策略。不同机器人命中落在同一深度区间时会产生近似误差，稀疏通信实现才体现显著带宽优势。该工作限制了直接共享的内容，仿真也未给出统计意义上的抗信息泄露保证。

\paragraph{FlyMark：在静态高斯的基本颜色内写入模型域水印}
2026.10~\cite{2610_11364}。FlyMark 将隐式所有权消息写入已经训练好的静态高斯资产，仅在色差受限的条件下调整零阶颜色。读出依赖合法密钥和对模型参数的访问；无需改变位置、尺度或训练额外解码器。水印所作用的是模型本身，并非单张渲染图。

原文按场景和密钥组合分别考察正常读出、错密钥和模型变化后的结果，并用固定相机进行画质对照。适用条件包括可读取静态参数数组以及作者定义的模型域变换。所测结果不能解释为对任意再训练、蒸馏或模型重建的普遍鲁棒性。

\paragraph{PrinGS：在高斯表示层处理人脸和文字隐私}
2026.10~\cite{2610_36168}。PrinGS 根据敏感区域的图像掩码和光栅化时的可见性定位相关高斯，再对其局部外观或位置进行匿名化。它既可在现有场景上操作，也可把匿名化约束加入训练；现有场景模式可用渲染视图上的掩码，因此可以在原始训练图像缺失时使用。

作者在脸部和文字实验中分别检查检测率、再识别相似度或 OCR 字符错误率。效果受掩码覆盖、可见性估计以及是否采到了关键视角影响，留出视图上的隐私指标不能代替证明所有视角都无法恢复信息。
